\section{HUNL Subgame Construction}
\label{app:hunl}

\paragraph{The two native subgames.}
The native \emph{river} subgame (3{,}000 Infosets) deals non-conflicting hand pairs over
a fixed board and plays one street of check\slash bet\slash fold\slash call with
exact hand evaluation.
The \emph{turn} subgame (83{,}040 Infosets) adds a 44-card river chance layer between two
betting streets, with $K{=}12$ hands per player, sunk-cost payoffs, fractional
pot-based bet sizing, and all-in handling. These fixed-board, fixed-range trees
capture the repeated re-solving structure of HUNL-style workloads
\cite{burch2014solving,ganzfried2015endgame}.

\paragraph{Background: subgame solving and online re-solving.}
Decomposition solves subgames independently given summaries of the rest of the game
\cite{burch2014solving}, with safety refinements \cite{brown2017safe,zhang2021subgame},
practical endgame solving \cite{ganzfried2015endgame,moravcik2016refining}, and
depth-limited variants that cap the lookahead with value estimates
\cite{brown2018depth,kovavrik2023value,kroer2020limited}. DeepStack's
continual re-solving \cite{moravvcik2017deepstack,vsustr2019monte,sustr2020sound} and
Libratus's endgame solver
\cite{brown2018superhuman} run exactly this workload online, under second-scale time
budgets. Our two native HUNL subgames reproduce the structure of this workload, while
\cref{sec:hunl} separates their reusable compilation cost from marginal solver time.

Both native subgames use exact 7-card hand evaluation ported from a production poker
engine and validated against brute-force enumeration.

\paragraph{River subgame.}
Fixed board (default \texttt{Ks Js Th 7d 2c}), pot 20, stacks 100. A chance root deals
each player one of $K{=}100$ candidate hands, uniformly over non-conflicting pairs
(9{,}161 deals). One betting street: check/bet/fold/call with bet fractions
$\{0.5, 1.0\}$ of the pot, no raises. Payoffs at showdown are $\pm(\text{half pot} +
\text{matched bets})$; folds forfeit the half-pot plus any called amount. The default
instance has 3{,}000 Infosets on 137{,}415 States.

\paragraph{Turn subgame.}
Board \texttt{Ks Js Th 7d} plus a full 44-card river chance layer. Each player holds
one of $K{=}12$ hands ($131$ non-conflicting deals at the default seed); betting
occurs on the turn and again on the river, with fraction-of-pot sizing
$b = \min(\text{stack} - \text{committed},\, f \cdot \text{pot})$ and all-in handling
(no river betting after all-in). Payoffs use \emph{sunk-cost} semantics: a turn fold
loses the half-pot; a river fold additionally loses the caller's committed chips;
showdown transfers the half-pot plus matched wagers. The default instance has
83{,}040 Infosets on 433{,}610 States, and tree depth yielding 8 fused execution
blocks per pass.

\paragraph{Analytic anchors.}
Two closed-form checks pin absolute correctness of payoffs, chance weighting, and the
evaluator simultaneously. (i)~River nuts anchor: if one player's range contains only
the best possible hand, the equilibrium value in exact arithmetic is the half-pot,
$+10$; the implemented game's chance discretization shifts this to an enumerable
$+10.000033$, which the solver and a brute-force best-response enumeration both
reproduce to $10^{-10}$.
(ii)~Quads anchor: on the board \texttt{7s 7d 7h Ks} with hole cards \texttt{7c 2d},
four-of-a-kind wins all 44 rivers, forcing root value $+10$; the solver converges to
it with exploitability below $10^{-3}$ after 1{,}000 iterations.

\paragraph{Best-response oracle.}
A recursive exact oracle, weighted by counterfactual reach and expanded deepest
first, validates the generic evaluator on these games to $10^{-9}$. The oracle itself is
validated by brute-force enumeration of all pure strategies on reduced instances.

\paragraph{External reconstruction of the river subgame.}
The river subgame can be reconstructed outside our codebase, providing an independent
instance for the
head-to-head timings of \cref{sec:results}. We express it as an OpenSpiel
\texttt{universal\_poker} ACPC instance \cite{lanctot2019openspiel}: fixed five-card
board, \texttt{handReaches} marking the $K{=}100$ candidate hands per player, and one
betting round with two bet sizes. Three details decide whether the two games coincide.
The chips-behind parameter is set to the largest bet in our abstraction (below the
nominal stack), so that \texttt{universal\_poker}'s all-in action coincides with its
pot-size bet and is deduplicated away; at the nominal stack the external game acquires a
third bet size and is strictly larger. The \texttt{handReaches} slot order is
lexicographic in $(\mathrm{lo}, \mathrm{hi})$ with $\mathrm{card} = 4\,\mathrm{rank} +
\mathrm{suit}$ over ranks $2{\ldots}\mathrm{A}$ and suits $s, h, d, c$, which differs
from GPU-CFR's own renderer in both rank and suit order. The opening action order is a
permutation of GPU-CFR's, so strategies are compared by action name and never by position.
Under this mapping the two constructions agree exactly on 137{,}416 nodes and 9{,}161
deals, and on the payoff at every terminal.

The mapping is verified at the level of the iterate as well as the tree: after three
alternating CFR$^+$ iterations an independent from-scratch reference solver over the
external representation and GPU-CFR over its own representation reach exploitability
$1.6107099050007263$ and $1.6107099050007272$ respectively. The reference solver is itself pinned
against OpenSpiel's own CFR and CFR$^+$ on Kuhn poker at 1, 10, and 100 iterations to
$10^{-9}$; OpenSpiel's solvers cannot be run on the river subgame directly, because
\texttt{universal\_poker} enumerates all $1{,}624{,}350$ deals with the excluded ones at
probability zero and their tabular policy does not admit card removal. Exploitability
for the external arm is therefore scored by our own exact tabular best response, the
same evaluator used by every other arm.

\paragraph{Scope of the external reconstruction.}
The turn subgame is deliberately not reconstructed in \texttt{universal\_poker}. Its
all-in deduplication would have to hold on the turn and again on the river, and a
single chips-behind parameter cannot satisfy both: any value that collapses all-in
into the pot-size bet on one street splits them on the other. Because a game that differs from the one we solve would confound the comparison,
the \texttt{universal\_poker} arm
covers the river subgame and the turn subgame is reported against our own arms.

External HUNL head-to-heads do not depend on that reconstruction at all. The Libratus
endgames are read from the same released endgame files by both our builder and
PokerRL, and both external CFR baselines solve our own trees directly from the export
of \cref{app:libratus}, so those comparisons are on a single shared game.
